\section*{Hoofdstuk \ref{ch:g8:the-atmosphere} --- De atmosfeer}

\begin{solution}{exo:g8:the-atmosphere:1}
(a) \omterm{def:g8:the-atmosphere:layers}{stratosfeer}; (b) \omterm{def:g8:the-atmosphere:layers}{troposfeer}; (c) \omterm{def:g8:the-atmosphere:layers}{mesosfeer}; (d) \omterm{def:g8:the-atmosphere:layers}{thermosfeer}.
\end{solution}

\begin{solution}{exo:g8:the-atmosphere:2}
(a) \qty{15}{\degreeCelsius}. (b) Ongeveer \qty{-56}{\degreeCelsius}.
(c) Ongeveer \qty{-3}{\degreeCelsius}. (d) Ongeveer \qty{-86}{\degreeCelsius}.
(e) Van \qty{11}{km} tot ongeveer \qty{20}{km} (en over een paar kilometer
rond \qty{50}{km} en net boven \qty{86}{km}).
\end{solution}

\begin{solution}{exo:g8:the-atmosphere:3}
(a) Niet waar: ze daalt in de \omterm{def:g8:the-atmosphere:layers}{troposfeer}, stijgt in de \omterm{def:g8:the-atmosphere:layers}{stratosfeer}, daalt
in de \omterm{def:g8:the-atmosphere:layers}{mesosfeer} en stijgt weer in de \omterm{def:g8:the-atmosphere:layers}{thermosfeer}. (b) Waar. (c) Niet waar:
ze bevindt zich in de \omterm{def:g8:the-atmosphere:layers}{stratosfeer}. (d) Waar.
\end{solution}

\begin{solution}{exo:g8:the-atmosphere:4}
(a) De \omterm{def:g8:the-atmosphere:layers}{troposfeer}. (b) De \omterm{def:g8:the-atmosphere:layers}{stratosfeer} (in de standaardatmosfeer ligt de
\omterm{def:g8:the-atmosphere:layers}{tropopauze} op \qty{11}{km}). (c) De \omterm{def:g8:the-atmosphere:layers}{stratosfeer}. (d) De \omterm{def:g8:the-atmosphere:layers}{mesosfeer}.
\end{solution}

\begin{solution}{exo:g8:the-atmosphere:5}
(a) $20 - 6.5 \times 2 = 20 - 13 = \qty{7}{\degreeCelsius}$. (b) $20 -
6.5 \times 3 = 20 - 19.5 = \qty{0.5}{\degreeCelsius}$. (c) De lucht moet
20 graden afkoelen: $20 \div 6.5 \approx 3.1$, dus op ongeveer
\qty{3.1}{km}.
\end{solution}

\begin{solution}{exo:g8:the-atmosphere:6}
(a) $3 - 0.5 = \qty{2.5}{km}$. (b) $6.5 \times 2.5 = 16.25$: ongeveer 16
graden. (c) $18 - 16.25 \approx \qty{2}{\degreeCelsius}$.
\end{solution}

\begin{solution}{exo:g8:the-atmosphere:7}
(a) \qty{500}{hPa}, \qty{250}{hPa} en \qty{125}{hPa}. (b) Dicht bij de
grond is de regel goed (\qty{500}{hPa} tegenover \qty{506}{hPa}). Hoger
geeft hij te veel: \qty{250}{hPa} tegenover \qty{227}{hPa}, en
\qty{125}{hPa} tegenover \qty{96}{hPa}. De halveringsregel is maar een
ruwe leidraad.
\end{solution}

\begin{solution}{exo:g8:the-atmosphere:8}
(a) $10~\text{hPa} = 10^1$~hPa; $10^3 \div 10^1 = 10^2$: honderd keer
kleiner. (b) $0.01~\text{hPa} = 10^{-2}$~hPa. (c) $10^3 \div 10^{-2} =
10^5$: honderdduizend keer kleiner.
\end{solution}

\begin{solution}{exo:g8:the-atmosphere:9}
(a) $6 \div 1.5 = 4$ keer breder. (b) Terwijl hij stijgt, drukt de lucht
eromheen steeds minder hard, dus zet het gas binnenin uit en groeit de
ballon tot zijn wand scheurt. (c) Nadat de ballon is geknapt, remt de
parachute de val van de radiosonde af, zodat die niemand verwondt en niets
beschadigt.
\end{solution}

\begin{solution}{exo:g8:the-atmosphere:10}
(a) $20 \div 7 \approx 2.9$ keer dikker. (b) Dichter bij die van de \omterm{def:g3:continents-and-oceans:poles}{polen}:
\qty{11}{km} ligt \qty{4}{km} boven \qty{7}{km} maar \qty{9}{km} onder
\qty{20}{km} (het profiel beschrijft de gematigde breedten, niet de
\omterm{def:g3:continents-and-oceans:equator}{evenaar}). (c) Aan het begin ligt de \omterm{def:g8:the-atmosphere:layers}{tropopauze} op ongeveer \qty{7}{km},
onder het vliegtuig: het vliegtuig bevindt zich in de \omterm{def:g8:the-atmosphere:layers}{stratosfeer}. Aan het
eind ligt de \omterm{def:g8:the-atmosphere:layers}{tropopauze} op ongeveer \qty{20}{km}: het vliegtuig bevindt
zich in de \omterm{def:g8:the-atmosphere:layers}{troposfeer}.
\end{solution}

\begin{solution}{exo:g8:the-atmosphere:11}
(a) $1 - 506 \div 1013 \approx 0.50$: \textbf{50~\%}. (b) $1 - 227 \div
1013 \approx 0.78$: \textbf{78~\%}. (c) $1 - 12 \div 1013 \approx 0.99$:
\textbf{99~\%}. (d) Ongeveer 1~\%.
\end{solution}

\begin{solution}{exo:g8:the-atmosphere:12}
(a) $6371 \div 15 \approx 425$~km per centimeter. (b) $100 \div 425
\approx 0.24$~cm, ongeveer \qty{2.4}{mm}. (c) $11 \div 425 \approx
0.026$~cm, ongeveer \qty{0.26}{mm}. (d) $0.26 \div 0.1 \approx 2.6$: twee
of drie vellen papier.
\end{solution}

\begin{solution}{pb:g8:the-atmosphere:1}
\textbf{1.} $300 \times 10 = 3000$~m.

\textbf{2.} $6000 \div 300 = 20$ minuten; $\num{30000} \div 300 = 100$
minuten.

\textbf{3.} 100 minuten, dat is 1~uur en 40~minuten.

\textbf{4.} Op \qty{8}{km}: de \omterm{def:g8:the-atmosphere:layers}{troposfeer} (de temperatuur daalt nog). Op
\qty{30}{km}: de \omterm{def:g8:the-atmosphere:layers}{stratosfeer}.

\textbf{5.} $15 - (-50) = 65$ graden.

\textbf{6.} $65 \div 10 = \qty{6.5}{\degreeCelsius}$ per kilometre.

\textbf{7.} Van \qty{12}{km} tot \qty{20}{km} blijft ze
\qty{-56}{\degreeCelsius}. De afkoeling stopt tussen \qty{10}{km} en
\qty{12}{km}: de \omterm{def:g8:the-atmosphere:layers}{tropopauze} ligt op ongeveer \qty{11}{km}.

\textbf{8.} Ze stijgt: \qty{-55}{\degreeCelsius} op \qty{22}{km},
\qty{-47}{\degreeCelsius} op \qty{30}{km}. De ballon bevindt zich in de
\omterm{def:g8:the-atmosphere:layers}{stratosfeer}, die wordt opgewarmd door het ozon dat het ultraviolette licht
van de Zon absorbeert.

\textbf{9.} $1013 \div 2 \approx \qty{507}{hPa}$; die wordt bereikt tussen
\qty{4}{km} (\qty{617}{hPa}) en \qty{6}{km} (\qty{472}{hPa}).

\textbf{10.} $1013 \div 4 \approx \qty{253}{hPa}$.

\textbf{11.} Tussen \qty{10}{km} (\qty{265}{hPa}) en \qty{12}{km}
(\qty{194}{hPa}).

\textbf{12.} Vanaf \qty{10}{km} moet de druk nog $265 - 253 =
12$~hPa \omterm{def:g1:landscapes-around-us:landform}{dalen}, van de 71~hPa die hij over 2~km verliest: $12 \div 71 \times 2
\approx 0.34$~km. De druk is een kwart van zijn waarde op \omterm{def:g4:world-physical-map:altitude}{zeeniveau} op
ongeveer \textbf{\qty{10.3}{km}}.
\end{solution}
